% Book pp. 125-128 (PDF pp. 126-129)
\section{Linear and Quadratic Factor Theorems}

The linear quadratic factor theorem is another method which can be used to solve
polynomial functions. This is when we break polynomial functions into linear and
quadratic factors to make solving it simpler.

For example, to use the linear and quadratic factor approach to find the factors of
$x^4 - 9x^2 + 14$, we:

\textbf{Step 1}: Let $x^2 = u$

\textbf{Step 2}: Using this substitution it means $x^4 - 9x^2 + 14 = (u)^2 - 9(u) + 14$

\textbf{Step 3}: Factorise $u^2 - 9u + 14$:

$u^2 - 7u - 2u + 14$

$u(u - 7) - 2(u - 7)$

$(u - 2)(u - 7)$

\textbf{Step 4}: Substitute $x^2$ back in for $u$:

$(x^2 - 2)(x^2 - 7)$

$x^2 - 2 = 0$ and $x^2 - 7 = 0$

For $x^2 - 2 = 0$

$x^2 = 2$

$x = \pm\sqrt{2}$

$x = -\sqrt{2}$ or $x = \sqrt{2}$

$x + \sqrt{2}$ and $x - \sqrt{2}$ are factors

For $x^2 - 7 = 0$

$x^2 = 7$

$x = \pm\sqrt{7}$

$x = -\sqrt{7}$ or $x = \sqrt{7}$

$x + \sqrt{7}$ and $x - \sqrt{7}$ are factors

Hence the factors are $+\sqrt{2}$, $x - \sqrt{2}$, $x + \sqrt{7}$ and $x - \sqrt{7}$

\begin{example}{Example 3.17}
Given $f(x) = x^4 - 1$, find all the factors using the linear quadratic factor approach.

\solution
Let $u = x^2$

$f(x) = (x^2)^2 - 1$

$(u)^2 - 1$

$u^2 - 1$

Factorising:

$(u - 1)(u + 1)$

Substituting $x^2$ back in:

$(x^2 - 1)(x^2 + 1)$

$(x^2 - 1)(x^2 + 1) = 0$

$x^2 - 1 = 0$ or $x^2 + 1 = 0$

For $x^2 - 1 = 0$

$x^2 = \pm 1$

$x = 1$ or $x = -1$

For $(x^2 + 1)$

$x^2 + 1 = 0$

$x^2 = -1$

$x = \pm\sqrt{-1}$

\textbf{NB:} $\sqrt{-1} = \boldsymbol{i}$

$x = \pm i$

$x = -i$ or $x = i$

Hence the factors are $(x + 1)$, $(x - 1)$, $(x - i)(x + i)$
\end{example}

\begin{example}{Example 3.18}
Given $f(x) = 2x^4 + x^2 - 10$, find all the factors using the linear quadratic factor
approach

\solution
Let $u = x^2$

$f(x) = 2(x^2)^2 + u - 10$

$2(u)^2 + u - 10$

$2u^2 + u - 10$

Factorising:

$2u^2 - 4u + 5u - 10$

$2u(u - 2) + 5(u - 2)$

$(2u + 5)(u - 2)$

Substituting:

$u = x^2$

$(2x^2 + 5)(x^2 - 2)$

$(2x^2 + 5)(x^2 - 2) = 0$

$2x^2 + 5 = 0$ or $x^2 - 2 = 0$

For $2x^2 + 5 = 0$

$x^2 = -\frac{5}{2}$

$x = \pm\sqrt{-\frac{5}{2}}$

$x = \frac{5}{2}\,i$ or $x = -\frac{5}{2}\,i$
\booknote[S3-15]{$\sqrt{-\frac{5}{2}} = \sqrt{\frac{5}{2}}\,i = \frac{\sqrt{10}}{2}\,i$,
not $\frac{5}{2}\,i$; the two complex factors in the last line inherit the slip.}

For $(x^2 - 2)$

$x^2 - 2 = 0$

$x^2 = 2$

$x = \pm\sqrt{2}$

$x = -\sqrt{2}$ or $x = \sqrt{2}$

Hence the factors are $\left(x + \frac{5}{2}\,i\right)\left(x - \frac{5}{2}\,i\right)
\left(x - \sqrt{2}\right)\left(x + \sqrt{2}\right)$
\end{example}

We can use what we have learned to help in determining the maximum or minimum
profit or productivity in the real world.

In small groups, or individually, discuss how the following problem can be solved.

\begin{example}{Example 3.19}
A water production company made sales of $x$ bags of water and the profit was
given by $2x^2 + 7x - 15$ where $0 \le x \le 10$.

Help the company determine their maximum profit and loss.

\solution
Since the range was given \textbf{(}$0 \le x \le 10$), we do our calculation within
this range.

Evaluating the function in the given interval, we obtain:
\begin{center}
\renewcommand{\arraystretch}{1.3}
\begin{tabular}{|l|c|c|c|c|c|c|c|c|c|c|c|}
\hline
x & 0 & 1 & 2 & 3 & 4 & 5 & 6 & 7 & 8 & 9 & 10 \\
\hline
$2x^2 + 7x - 15$ & -15 & -6 & 7 & 24 & 45 & 70 & 99 & 132 & 169 & 210 & 255 \\
\hline
\end{tabular}
\end{center}
The maximum profit occurs when x=10 with a profit of 255 while the loss will
occur when there is no production, i.e., $x = 0$ at a loss of 15 ( -15).
\end{example}

\section*{Extended Reading}
\begin{itemize}
  \item Lial, M. L., Hornsby, E. J., \& McGinnis, T. (2012). Algebra for college
  students. (7th Ed. Pearson Education, Inc). Pages 346 -- 375.
\end{itemize}
